\textbf{Momento angular, producto vectorial y sistema de ecuaciones.} Primero se calculan las tres filas de $\boldsymbol H=I\boldsymbol\omega$. Las entradas $xz$ y $zx$ son negativas con el convenio del enunciado.\[\fitmath{\boldsymbol H=\left[\begin{matrix}1168.0\\990.0\\-1765.0\end{matrix}\right]\,\mathrm{kg\,m^2/s}}\]\[\fitmath{H_{1}=(11000)(0.1)+(-850)(-0.08)\approx1168\,\mathrm{kg\,m^2/s}}\]\[\fitmath{H_{2}=(16500)(0.06)\approx990\,\mathrm{kg\,m^2/s}}\]\[\fitmath{H_{3}=(-850)(0.1)+(21000)(-0.08)\approx-1765\,\mathrm{kg\,m^2/s}}\]\[\fitmath{c_{1}=(0.06)(-1765)-(-0.08)(990)\approx-26.7\,\mathrm{N\,m}}\]\[\fitmath{c_{2}=(-0.08)(1168)-(0.1)(-1765)\approx83.06\,\mathrm{N\,m}}\]\[\fitmath{c_{3}=(0.1)(990)-(0.06)(1168)\approx28.92\,\mathrm{N\,m}}\]El producto giroscópico se desarrolla como $(qH_z-rH_y,\ rH_x-pH_z,\ pH_y-qH_x)$; su valor es:\[\fitmath{\boldsymbol\omega\times\boldsymbol H=\left[\begin{matrix}-26.7\\83.06\\28.92\end{matrix}\right]\,\mathrm{N\,m}}\]Se resta de los momentos aplicados, obteniendo el segundo miembro $\boldsymbol b=\boldsymbol M-\boldsymbol\omega\times\boldsymbol H$.\[\fitmath{\boldsymbol b=\left[\begin{matrix}726.7\\-533.06\\351.08\end{matrix}\right]\,\mathrm{N\,m}}\]La ecuación de $y$ queda desacoplada: $\dot q=b_y/I_{yy}$. Para $x,z$, eliminar la otra aceleración en el sistema $I_{xx}\dot p-I_{xz}\dot r=b_x$, $-I_{xz}\dot p+I_{zz}\dot r=b_z$ da:\[\fitmath{\dot p=\frac{I_{zz}b_x+I_{xz}b_z}{I_{xx}I_{zz}-I_{xz}^2},\qquad\dot r=\frac{I_{xz}b_x+I_{xx}b_z}{I_{xx}I_{zz}-I_{xz}^2}}\]\[\fitmath{(\dot p,\dot q,\dot r)=\left[\begin{matrix}0.0675668182953176 & -0.0323066666666667 & 0.0194529426452867\end{matrix}\right]\,\mathrm{s^{-2}}}\]Para comprobar, multiplicar el vector de aceleraciones por $I$ y sumar el término giroscópico recupera $\boldsymbol M$. El trabajo giroscópico se anula por perpendicularidad; la tasa de energía es $\boldsymbol\omega\cdot\boldsymbol M=12.6\,\mathrm W$.
